\chapter{Einleitung}
\label{chap:einleitung}

% Inhalt laut Vorgabe:
%  - Beschreibung der gesamten Anlage (Kurzfassung)
%  - Beschreibung des eigenen Moduls (detailliert)

\section{Gesamtanlage}
\label{sec:gesamtanlage}

Die Modellanlage (\abkref{VPS}) bildet die Herstellung von Popcorn in vier
hintereinander geschalteten Modulen ab. Der Mais durchläuft die Module in der
Reihenfolge Delivery, Storage, Dosing und Production. Die Anlage wird
ausschließlich über einen \abkref{MQTT}-Broker angesprochen: Sensorwerte werden
vom Broker gelesen, Aktoren werden durch Nachrichten an den Broker geschaltet
\cite{doku-popcorn}.

\begin{table}[htbp]
  \centering
  \begin{tabularx}{\textwidth}{@{}l l X@{}}
    \toprule
    \textbf{Modul} & \textbf{Name} & \textbf{Aufgabe} \\
    \midrule
    1 & Delivery   & Anlieferung: Mais aus dem Silo per Förderband und Druckluft an Modul~2 übergeben \\
    2 & Storage    & Zwischenlagerung; Weitertransport über pneumatischen Muskel und Druckluft \\
    3 & Dosing     & Dosierung über Förderschnecke (Drehzahl per \abkref{PWM}) und Wägezelle \\
    4 & Production & Poppen mit Heizspirale; Becher und Flanke für die Ausgabe \\
    \bottomrule
  \end{tabularx}
  \caption{Module der Gesamtanlage}
  \label{tab:module}
\end{table}

\todo{Kurzbeschreibung der Gesamtanlage ergänzen: Zusammenspiel der Module,
welche Gruppe welches Modul steuert, Material- und Signalfluss.}

\section{Modul 1 -- Delivery}
\label{sec:modul-delivery}

Modul~1 ist das erste Modul der Anlage. Es fördert den Mais aus seinem Silo über
ein Förderband in einen Trichter. Von dort wird der Mais mit Druckluft durch einen
Schlauch zu Modul~2 (Storage) transportiert. Sobald der Motor des Förderbands
eingeschaltet wird, läuft die Druckluft automatisch mit \cite{doku-popcorn}.

Das Modul besitzt fünf Aktoren (Motor, Druckluft und drei Signallampen) sowie vier
Sensoren (Silo oben, Silo unten, Trichter und Not-Aus). Eine vollständige Liste mit
Funktion und MQTT-Topic steht in Kapitel~\ref{chap:hardware}.

\bildoderplatzhalter[width=0.8\linewidth]{abbildungen/delivery-schema.png}
  {Aufbau von Modul~1 (Delivery) mit Lage der Sensoren und Aktoren \cite{doku-popcorn}}
  {fig:delivery-schema}

Die Steuerung ist in Python mit der Bibliothek \texttt{paho-mqtt} umgesetzt
\cite{paho}. Sie besteht aus drei Schichten:
\begin{itemize}
  \item \textbf{Komponenten} (\texttt{komponenten.py}): kapseln je einen Aktor
        und senden dessen Befehle an den Broker.
  \item \textbf{Modul} (\texttt{Module/delivery.py}): hält die aktuellen
        Sensorwerte und verteilt eingehende Nachrichten.
  \item \textbf{Zustandsautomat} (\texttt{zustandsautomat.py}): entscheidet
        zyklisch alle 0,1~s anhand der Sensorwerte über den nächsten
        Zustand und schaltet die Aktoren über das Modul.
\end{itemize}

\todo{Aufgabe und Randbedingungen des Moduls detailliert beschreiben
(Startbedingung, Abschaltbedingungen, Kopplung zu Modul~2, Sicherheitsverhalten).}
